\section{Data}\label{data}
\subsection{Sources}
We utilize a rich administrative dataset of public procurement auctions conducted via the BEC platform in São Paulo state. The dataset covers the period from January 2009 to December 2019 and includes detailed bid-level information for auctions of common goods and services. During this period, 1,344 PBUs (encompassing state-level executive agencies, the legislative and judicial branches, as well as some municipal governments and even a few private organizations authorized to use the system) carried out procurement through BEC. In total, the data comprise 3,866,407 individual bidding transactions submitted by 19,007 distinct firms, corresponding to 169,607 unique procurement items. Each record in the data identifies the purchase offer (PO) and item, the participating firm, the bid amount, and the auction outcome (win or loss), among other details.

The BEC system's hierarchical catalog organizes items into groups and classes. For example, medical supplies are grouped under a specific category (e.g., group 65 for medical and hospital items), with finer classifications down to individual item codes.\footnote{An illustrative item is code 110639, which corresponds to "Furosemide 40mg tablets," classified under class 6531 (a category of pharmaceuticals) within group 65.} This structured coding is useful for analyzing procurement patterns by item type or sector. Our analysis leverages the full breadth of goods and services in the dataset, while controlling for any category-specific effects as needed.

For the purposes of our study, we focus on competitive tenders that resulted in a contract award (i.e., at least one winner and one or more losing bids). We exclude tenders that were canceled or received only a single bid, since those do not provide information on competitive dynamics. This yields a sample of auctions where a meaningful competition took place, which is essential for implementing our regression discontinuity approach. 

\subsection{Variable definitions}
We construct key variables to capture firm outcomes and auction characteristics. First, we define a market such that two firms belong to the same market if they participated at least once in the same auction for the same type of good. This market classification is transitive, resulting in 7,581 distinct markets defined within at least one of our defined time periods.

Following \citet{kawai2022study}, we construct an \textit{Incumbent} dummy variable based on this market definition. Specifically, this variable equals 1 if a given firm won the last auction it participated in within the same market before the current auction, and 0 otherwise. This indicator reflects whether the firm entered the current tender as a reigning winner (incumbent) in the specific market or not.

Additionally, we define the \textit{Last Bid} dummy variable equal to 1 if the firm was the last to submit a bid in the auction, and 0 otherwise. %Second, we measure each firm's \textit{Backlog} as the total value of contracts the firm has won in the recent past. In practice, we calculate the backlog as the sum of all contract awards obtained by the firm in the 90 days (approximately three months) prior to a given auction. Because firms differ greatly in size and bidding frequency, we standardize this backlog measure by subtracting the firm's mean backlog (over the sample period) and dividing by its standard deviation. This standardization (also used by \citet{kawai2022study}) yields a normalized index of how "busy" a firm is relative to its own typical level of contract work at that point in time.

Finally, we define the \textit{Margin of Victory}, which serves as the running variable in our RDD design. For each auction, the margin of victory is calculated as the percentage difference between the winning bid and the second-lowest bid, expressed as a fraction of the winning bid. Formally, if $b_{\text{win}}$ is the winning bid and $b_{\text{runner-up}}$ is the next lowest bid, the margin is $(b_{\text{runner-up}} - b_{\text{win}})/b_{\text{win}}$. When this margin is very small, the auction outcome was decided by a tiny price difference.
